\chapter{Mind Your Motive}

These notes follow Michael Sandel's presentation of Kant as the point in the course where morality and freedom become inseparable. The lecture is organized with unusual care: first we are told why Kant matters, then we are shown why freedom cannot mean mere desire-satisfaction, then we are asked what gives an action moral worth, and only after that are we brought to the categorical imperative itself. The present chapter keeps that order, while preserving the examples and objections that make the argument move. This chapter is prepared for the LazyLearn track, with curation by LazyingArt LLC.

\section{Why Kant, and Why Now}

Sandel opens by telling us that Kant is the hardest philosopher in the course, but also one of the most important. The difficulty is justified by the scope of the project. Kant's moral philosophy is not presented as an isolated doctrine about a few duties; it is presented as an attempt to answer two large questions at once: what the supreme principle of morality is, and how freedom is possible.

The first decisive move is anti-utilitarian. Kant does not deny that we seek pleasure and avoid pain. In that limited sense, the utilitarians are, as Sandel puts it, half right. But they are wrong to treat pain and pleasure as our sovereign masters. What is distinctive about us is not merely that we can suffer or enjoy, but that we are rational beings capable of acting freely.

The lecture compresses that claim into a simple structural dependence:
\begin{equation}
\text{rational capacity} \Longrightarrow \text{dignity and respect}.
\end{equation}
This is why the person is not sacred because he owns himself, and not sacred because his happiness enters a social calculus, but because he is a bearer of reason. The anti-utilitarian conclusion follows immediately: if rational agency is the source of dignity, then persons may not be used merely as instruments, even for good ends.

\section{Rational Beings, Dignity, and Freedom Beyond Desire}

From here Sandel turns to freedom. He begins with the ordinary understanding and then rejects it. Freedom is often taken to mean doing what we want, or encountering no obstacle to the satisfaction of our wants. Kant's point is that this account is too weak, because it leaves the will under the rule of desires it did not choose.

We can state the contrast in the lecture's own schematic form:
\begin{align}
\text{freedom} &\neq \text{doing what we want},\\
\text{freedom} &= \text{autonomy}.
\end{align}
And autonomy is immediately opposed to heteronomy:
\begin{align}
\text{autonomy} &= \text{acting according to a law I give myself},\\
\text{heteronomy} &= \text{acting according to an inclination or desire not self-given}.
\end{align}

The lecture then sharpens the point by way of causal language. If I act under appetite, I am not authoring the law of my action. I am obeying a prompt I find already there. This is why Sandel treats the Sprite slogan, ``Obey your thirst,'' as a nearly perfect Kantian slogan. It reveals precisely what is wrong with identifying free choice with preference-satisfaction: the command comes from thirst, not from reason.

The chain of dependence is therefore:
\begin{equation}
\text{inclination} \to \text{heteronomy} \to \text{necessity},
\qquad
\text{reason} \to \text{autonomy} \to \text{freedom}.
\end{equation}
This is also why the billiard-ball example matters. A billiard ball falls according to gravity; we do not call that free. Kant's claim is that when our will is moved in the same merely causal way by appetite and impulse, we are not yet free in the morally relevant sense.

\subsection{Question \& Answer}

\paragraph{Question.} If satisfying my desires feels free, why does Kant say it is a form of unfreedom?

\paragraph{Answer.} Because the feeling of spontaneity is not yet authorship. I may choose Sprite rather than Pepsi, but if what governs the choice is thirst, or even a desire shaped by advertising, then the governing principle of the action is not mine in the relevant sense. Kant's notion of freedom is more demanding: freedom is the power to act from a law we legislate through reason, not merely the power to execute a desire efficiently.

This is why Sandel can move from freedom to dignity without changing the subject. If autonomy means that we can act from laws we give ourselves, then we are not merely bundles of appetite. We are beings who can stand in judgment on appetite.

\section{Autonomy, Morality, and the Motive of Duty}

Once freedom has been redefined, morality must be redefined with it. Sandel's pivot is exact: if freedom is autonomy, then morality cannot be about the successful production of outcomes alone. It must concern the quality of the will from which the action proceeds.

The lecture's moral compression is:
\begin{equation}
\text{moral worth} \Longrightarrow \text{depends on motive, not consequences}.
\end{equation}
Kant's famous formulation about the good will appears here with full force. A good will is good not because of what it accomplishes, but because it is good in itself:
\begin{equation}
\text{good will} \Longrightarrow \text{good in itself}.
\end{equation}
And the central sufficiency claim is:
\begin{equation}
\text{right action} + \text{duty as motive} \Longrightarrow \text{moral worth}.
\end{equation}
Conversely,
\begin{equation}
\text{right action} + \text{self-interest alone} \not\Longrightarrow \text{moral worth}.
\end{equation}

The lecture then proceeds experimentally through cases designed to isolate motive.

\paragraph{The shopkeeper.} The shopkeeper gives correct change, but only because cheating would damage his reputation and cost him business. The act conforms to duty, but it is not done from duty. This is Sandel's cleanest illustration of the gap between outward rightness and moral worth.

\paragraph{The duty against suicide.} Here the point is harder to see in ordinary life, because most people have many reasons for continuing to live. So Kant isolates the relevant variable by imagining a miserable person with no inclination toward life who nevertheless preserves himself from duty. Again the structure is not about outward result but about motive.

\paragraph{The Better Business Bureau and the honor-code discount.} These cases widen the field without changing the point. ``Honesty is the best policy'' is precisely not yet a moral claim in Kant's sense if what makes honesty attractive is profit. And an honor pledge sustained by discounts looks like the same shopkeeper structure in another setting.

Sandel summarizes the connection between morality and freedom by saying that only when we act from duty, resisting inclination and self-interest, are we acting freely in the demanding Kantian sense. The lecture's underlying claim may be written:
\begin{equation}
\text{act from duty} \Longrightarrow \text{autonomy},
\qquad
\text{act from inclination} \Longrightarrow \text{heteronomy}.
\end{equation}

\subsection{Question \& Answer}

\paragraph{Question.} Can an action still have moral worth if duty is accompanied by self-interest, sentiment, or even the desire to be moral?

\paragraph{Answer.} Sandel's answer is subtle. If self-interest is the ground of the act, moral worth is lost. But the presence of supporting sentiments does not by itself destroy moral worth. The lecture introduces reverence for the moral law here: the will may be formed so as to honor the law, and feelings may support the action, provided they are not what justifies it. In compact form,
\begin{equation}
\text{supporting sentiments} \neq \text{ground of moral worth}.
\end{equation}
This is the point at which the apparently austere Kantian picture becomes more precise. The issue is not emotional sterility; it is what finally explains why the act was done.

\section{One Moral Law or Many?}

The next objection is deeper. If autonomy means acting according to a law I give myself, why is morality not subjective? Why does self-legislation not fragment into as many moral laws as there are persons in the room?

Sandel treats this as a genuine pressure on the theory. If the moral law were merely my private conscience, autonomy would collapse into idiosyncrasy. Kant's answer is that the reason that legislates morally is not private reason in the biographical sense. It is shared practical reason, the same rational capacity present in all persons as rational beings.

The lecture's answer is therefore:
\begin{equation}
\text{shared practical reason} \Longrightarrow \text{one moral law, not many private laws}.
\end{equation}
This is why autonomy is not self-expression. The law is self-given, but the self that gives it is the rational will as such, not a contingent bundle of interests, upbringing, or temperament. Sandel emphasizes this by describing practical reason as a priori with respect to particular empirical ends. Our interests vary; the rational capacity that legislates the moral law does not.

\subsection{Question \& Answer}

\paragraph{Question.} If the moral law is self-given, why is morality not merely subjective?

\paragraph{Answer.} Because the relevant self is not the empirical personality with its contingent desires and history. The legislating source is the universal rational capacity that persons share. Kant's autonomy is therefore not ``I make up my own morality,'' but ``the rational will gives itself the law that any rational being, as rational, must be able to acknowledge.'' The lecture thus moves from self-legislation to universality rather than from self-legislation to relativism.

\section{The Internal Reset and the Three Contrasts}

At this point Sandel performs an explicit reset. He turns back to Kant and says, in effect, that the groundwork is about two big questions: the supreme principle of morality and the possibility of freedom. Then he says that one way through the book is to keep three related contrasts constantly in view.

The three contrasts can be written in one line:
\begin{equation}
\text{duty} \leftrightarrow \text{inclination},
\qquad
\text{autonomy} \leftrightarrow \text{heteronomy},
\qquad
\text{categorical imperative} \leftrightarrow \text{hypothetical imperative}.
\end{equation}
This reset is not repetition. It reorganizes the lecture into a formal structure.

Before moving to the third contrast, Sandel briefly returns to the first with the spelling-bee example. The student's confession is morally interesting precisely because the motive is mixed. Judith's defense of Kant is that the presence of an additional motive such as not wanting to feel like a slime does not destroy moral worth provided duty remains the reason for acting. This closes the earlier discussion of mixed motives and lets Sandel move forward.

He then states the second contrast again: the will can be determined autonomously or heteronomously. Only in the first case is it free. And from there he turns to the third contrast: the two kinds of commands reason can issue.

A hypothetical imperative has the familiar instrumental form
\begin{equation}
\text{If you want } X,\ \text{do } Y.
\end{equation}
The shopkeeper provides the lecture's clean example: if you want a good business reputation, do not shortchange customers.

By contrast, the categorical imperative is not ordered to a further purpose. It commands unconditionally. Sandel's slide marks the contrast exactly, and it should be kept as evidence of the lecture's structure.

\begin{figure}[htbp]
\centering
\includegraphics[width=0.82\textwidth]{lecture_06_figure_02.png}
\caption{Kant on hypothetical and categorical imperatives.}
\end{figure}

The cleaned reconstruction that belongs beside the slide is:
\begin{align}
\text{good merely as means} &\Longrightarrow \text{hypothetical imperative},\\
\text{good in itself} &\Longrightarrow \text{categorical imperative}.
\end{align}
This lets Sandel connect the three contrasts in the strongest possible way. If morality requires duty rather than inclination, and if freedom requires autonomy rather than heteronomy, then the appropriate command of reason must be categorical rather than hypothetical.

\section{The Categorical Imperative as Universal Law}

Only now does Sandel ask what the moral law actually says. Kant offers several formulations, but the lecture first takes up the formula of universal law. A maxim is defined as the rule or principle that explains why we are acting, and the test is whether that maxim can be willed as universal law.

The formulation may be stated as:
\begin{equation}
\text{Act only on that maxim whereby you can at the same time will that it should become a universal law.}
\end{equation}

\paragraph{Worked example: false promise.} Sandel's worked example is the false promise. The case is simple: I need money, I know I cannot repay it soon, and I nevertheless promise repayment in order to secure the loan. The maxim can be written schematically as
\begin{equation}
\text{When I need money, I will promise repayment though I know I cannot repay.}
\end{equation}
The test proceeds step by step:
\begin{enumerate}
\item Formulate the maxim that explains the intended action.
\item Universalize it: suppose everyone acted on that maxim whenever in need.
\item Ask whether the practice presupposed by the maxim could survive its own universalization.
\item In this case, universalized false promising destroys the institution of promising.
\item Therefore the maxim defeats itself and cannot be willed as universal law.
\end{enumerate}
The lecture's compressed contradiction structure is
\begin{equation}
\text{false-promise maxim universalized} \Longrightarrow \text{promise practice collapses} \Longrightarrow \text{contradiction}.
\end{equation}

\subsection{Question \& Answer}

\paragraph{Question.} Does this test secretly rely on consequences after all?

\paragraph{Answer.} Sandel pauses on exactly this objection, because it sounds at first as if Kant is saying, ``Do not lie because the world will go badly if everybody lies.'' That would indeed be consequentialist reasoning. His answer is that universalization is a test, not the underlying reason. The point is not to calculate bad social outcomes. The point is to determine whether my maxim illicitly privileges my special circumstances over everyone else's, and whether it can be coherently adopted as a law for rational agents as such. The contradiction is not a utility calculation; it is a failure of universality.

\section{Humanity as an End}

Sandel then turns to the second formulation of the categorical imperative because it is, as he says, more intuitively accessible. The thought is that a categorical imperative cannot be grounded in someone's merely private ends. Its ground must lie in something with absolute value. Kant's answer is that rational beings exist as ends in themselves.

The formula is therefore:
\begin{equation}
\text{Treat humanity, whether in yourself or another, never simply as a means but always at the same time as an end.}
\end{equation}
Here the false promise becomes immediately transparent. In deceiving another person for the sake of money, I am manipulating his rational agency and subordinating it to my private purpose. I am not merely interacting with him; I am using him as material for my end.

The same structure is then applied to murder and suicide:
\begin{equation}
\text{murder},\ \text{suicide} \Longrightarrow \text{use humanity merely as means} \Longrightarrow \text{violate dignity}.
\end{equation}
This is why Sandel says that, for Kant, suicide is morally on a par with murder. In both cases what is violated is humanity itself, the rational nature that commands respect. And because that rational nature is universal rather than particular, Kantian respect differs from love, sympathy, altruism, or solidarity. Those may be directed toward particular people because of who they are. Respect, by contrast, is respect for rational being as such.

\subsection{Question \& Answer}

\paragraph{Question.} How can we use other people as means in everyday life without treating them merely as means?

\paragraph{Answer.} Patrick's classroom objection is exactly right to raise here, because everyday life is full of ordinary purposive interaction. We hire, buy, sell, ask directions, seek help, and cooperate in projects. Kant's prohibition is not against all such use. The distinction is:
\begin{equation}
\text{use others as means} \neq \text{use others merely as means}.
\end{equation}
Ordinary cooperation is compatible with the categorical imperative when it proceeds under conditions of respect for the other as a rational being. Fraud, coercion, manipulation, and degradation are objectionable because they collapse that distinction. The other person is then not a partner in action but only an instrument of my purpose.

\section{Summary}

The lecture has a strict order, and the order matters. We begin with Kant as the philosopher of dignity, then discover that dignity rests on rational agency. From there freedom is redefined as autonomy rather than desire-satisfaction. Once freedom is understood that way, morality can no longer be based on outcomes alone; it must depend on motive, and specifically on duty. The distinction between duty and inclination then leads to the distinction between autonomy and heteronomy, and finally to the distinction between hypothetical and categorical imperatives.

The two principal formulations of the moral law then fall into place. The formula of universal law tests whether a maxim can be willed universally. The formula of humanity as an end states more directly what is at stake: persons, as rational beings, must never be treated merely as means. The chapter therefore closes where the lecture closes: with the question still alive, not whether Kant has given us an easy morality, but whether he has given us the right one.